\documentclass[11pt]{article}
\usepackage[a4paper,margin=25mm]{geometry}
\usepackage{amsmath,amssymb}
\usepackage[T1]{fontenc}
\usepackage{lmodern}
\usepackage{hyperref}
\setlength{\parindent}{0pt}
\setlength{\parskip}{7pt}
\title{The filed injective harmonious labeling definition fails on a tree}
\author{SucRunBug}
\date{2 October 2026}
\begin{document}
\maketitle
\textbf{Conjecture 00000000251 -- FALSE.}

\textbf{Filed claim addressed.} Every tree has an injective vertex labeling modulo its edge count $m$, with edge sums bijective modulo $m$.

\section*{Proof}
The one-edge tree $K_2$ has two vertices and $m=1$. The residue ring $\mathbb Z/1\mathbb Z$ has exactly one element. Every vertex-labeling map to that ring is constant and cannot be injective. Thus no labeling satisfies the filed definition, even before checking edge sums.

The same cardinality obstruction applies to any nonempty finite tree: it has $m+1$ vertices and only $m$ labels. The submission uses $K_2$ to give a fully formal concrete instance.

In the usual harmonious-labeling formulation for trees, one permits a repeated vertex label. The filed definition instead requires injectivity on all vertices. The argument does not settle the standard version with that repetition allowed.

\section*{Formalization scope}
Lean proves K2 is an actual one-edge simple-graph tree and proves there is no injective map from Fin 2 to ZMod 1.

\textbf{Reproduce.} In \texttt{lean4/}, run \texttt{lake exe cache get},
\texttt{lake build}, then \texttt{lake env lean Check.lean}. Versions:
Lean/Mathlib \texttt{v4.33.0}. Independent finite checks are in
\texttt{reproduce.py}; they are not a substitute for the complete proof.

\textbf{Source.} \url{https://github.com/The-Last-Math-Competition/The-Last-Math-Competition/blob/28a22efac52c4a7abd4bba8a143b1c1a06b7e177/conjectures/00000000251.md}
\end{document}
